% 张三慧第四版第二卷第18章《磁力》独立候选解析
% 主教材PDF物理页157–161，图18.20–18.33需源图复审。
\section*{第18章\quad 磁力：独立解析}

\paragraph{习题 18.1}
保持高度不变要求竖直方向总力为零。
电量 \(q>0\) 的粒子水平运动，须
\(\boldsymbol F_B=q\boldsymbol v\times\boldsymbol B\)
\emph{竖直向上}且 \(F_B=mg\)。
取 \(\boldsymbol B\perp\boldsymbol v\)，
\[
 \boxed{B=\frac{mg}{qv}
 =\frac{(5.0\times10^{-4})(9.8)}
 {(2.5\times10^{-8})(6.0\times10^4)}
 \approx3.27\,\mathrm T}.
\]
例如速度沿 \(+\hat x\)，
场应沿 \(+\hat y\)，
使 \(\hat x\times\hat y=+\hat z\)。
一般只要求与 \(v\) 垂直的场分量
\(B_\perp=mg/(qv)\)；
沿运动方向另加场分量不改变竖直磁力。

\paragraph{习题 18.2}
电子以 \(v_0=10^7\,\mathrm{m/s}\)
在图18.20 A点竖直向上进入圆轨道，
并向右弯曲，故带负电电子的
\(q\boldsymbol v\times\boldsymbol B\)
必须向右，
磁场方向\emph{入纸面}。
圆直径10cm，轨道半径
\(R=0.050\,\mathrm m\)，
\[
 \boxed{B=\frac{m_ev_0}{eR}
 \approx1.14\times10^{-3}\,\mathrm T}.
\]
速度大小在纯磁场中不变，
半圆弧长 \(\pi R\)，
\[
 \boxed{t=\frac{\pi R}{v_0}
 \approx1.57\times10^{-8}\,\mathrm s}.
\]

\paragraph{习题 18.3}
质子初动能 \(K=2000\,\mathrm{eV}\)
远低于静能，
可用 \(v=\sqrt{2K/m_p}
 \approx6.19\times10^5\,\mathrm{m/s}\)。
垂直和平行磁场速度
\(v_\perp=v\sin89^\circ\)，
\(v_\parallel=v\cos89^\circ\)。
回旋角频率
\(\omega=eB/m_p\)，
所以
\[
 \boxed{T=\frac{2\pi m_p}{eB}
 \approx6.56\times10^{-7}\,\mathrm s},
\]
\[
 \boxed{r=\frac{m_pv\sin89^\circ}{eB}
 \approx6.46\times10^{-2}\,\mathrm m},
\quad
 \boxed{h=v\cos89^\circ T
 \approx7.09\times10^{-3}\,\mathrm m}.
\]
若质子速度方向反向，螺旋绕行方向按
正电荷的右手法则确定。

\paragraph{习题 18.4}
忽略相对论修正，
被 \(U=20000\,\mathrm V\) 加速的电子
\[
 v=\sqrt{\frac{2eU}{m_e}}
 \approx8.39\times10^7\,\mathrm{m/s}.
\]
磁场垂直速度，圆轨道半径
\[
 R=\frac{m_ev}{eB}\approx9.54\,\mathrm m.
\]
飞行距离 \(L=0.20\,\mathrm m\) 与半径比
\(L/R\ll1\)，偏转的横向距离
\[
 \boxed{y\simeq\frac{L^2}{2R}
 \approx2.10\times10^{-3}\,\mathrm m
 =2.10\,\mathrm{mm}}.
\]
对精细显像管而言毫米级漂移可见，
实际管内加速与偏转场并非完全均匀，
此为初级估算。

\paragraph{习题 18.5}
周长 \(L=240\,\mathrm m\)
给轨道半径 \(R=L/(2\pi)\)。
磁力使动量方向转弯，
\(eBv=pv/R\)，与动量是否相对论无关：
\[
 \boxed{B=\frac{p}{eR}
 =\frac{2\pi p}{eL}
 =\frac{2\pi(1.49\times10^{-18})}
 {(1.6022\times10^{-19})(240)}
 \approx0.243\,\mathrm T}.
\]
不可用静质量的非相对论动量
重新倒算给定的 \(p\)。

\paragraph{习题 18.6}
带电粒子的回转半径一般为
\(r=p_\perp/(eB)\)。
按题给动量全为垂直场分量
\(p_\perp=10^{-16}\,\mathrm{kg\,m/s}\)：
\[
 \boxed{r(B=10^{-8}\,\mathrm T)
 =\frac{10^{-16}}{e\,10^{-8}}
 \approx6.24\times10^{10}\,\mathrm m},
\]
\[
 \boxed{r(B=10^8\,\mathrm T)
 \approx6.24\times10^{-6}\,\mathrm m}.
\]
两处磁场相差 \(10^{16}\) 倍，
半径因而反比变化；
电子动量很大，必须用相对论动量而非
\(m_ev\) 套非相对论数值。

\paragraph{习题 18.7}
\(B=20\,\mathrm T\)，
质子轨迹半径 \(R=3.5\,\mathrm m\)，
\[
 \boxed{p=eBR
 \approx1.12\times10^{-17}\,\mathrm{kg\,m/s}}.
\]
由相对论质心能量关系
\(E^2=(pc)^2+(m_pc^2)^2\)：
\[
 \boxed{E\approx21.01\,\mathrm{GeV}},
\quad
 \boxed{K=E-m_pc^2
 \approx20.07\,\mathrm{GeV}}.
\]
题若“能量”指总能量应报告 \(E\)；
若要求动能需扣除质子静能约0.938GeV。
\textbf{不可用} \(\tfrac12m_pv^2\)
估算20GeV量级的高能质子。

\paragraph{习题 18.8}
在地磁场中电子的回旋半径
\(r=m_ev_\perp/(eB)\)。
按原题的\emph{垂直速率保持
\(v_\perp=0.80\times10^8\,\mathrm{m/s}\)}
的估算：
\[
 \boxed{r_1=\frac{m_e(0.8\times10^8)}
 {e(4\times10^{-7})}
 \approx1.14\times10^3\,\mathrm m},
\]
\[
 \boxed{r_2=\frac{m_e(0.8\times10^8)}
 {e(2\times10^{-5})}
 \approx22.7\,\mathrm m}.
\]
这里作的是同一垂直速度的
两个局部场强估算；实际沿
梯度磁场缓慢运动，
磁矩绝热不变量可能导致
\(v_\perp^2/B\) 近似不变，
从而回转半径不严格按 \(1/B\)
而是按 \(1/\sqrt B\) 缩放。
题目没有给入射螺旋角，不足以进行
完整磁镜运动轨迹模拟。

\paragraph{习题 18.9}
氘核质量
\(m_D=3.34\times10^{-27}\,\mathrm{kg}\)，
电量 \(e\)，磁场 \(B=1.5\,\mathrm T\)。
非相对论回旋频率
\[
 \boxed{f=\frac{eB}{2\pi m_D}
 \approx1.15\times10^7\,\mathrm{Hz}
 =11.5\,\mathrm{MHz}}.
\]
D盒直径75cm，所以
可达到的最大半径 \(R_{\max}=.375\,\mathrm m\)，
出射最大动能
\[
 \boxed{K_{\max}=
 \frac{(eBR_{\max})^2}{2m_D}
 \approx7.59\,\mathrm{MeV}}.
\]
\emph{高频电源电压幅度}并不能直接代替最终
粒子能量；本题由轨道半径限制出射动量。

\paragraph{习题 18.10}
速度选择器中直线穿行的离子满足
\(qE=qvB\)，因此
\(\boxed{v=E/B}\)。
在分析场 \(B'\) 中
圆轨道半径
\(R=mv/(qB')\)，
由图18.22的半圆偏转
\(\ell=2R\)，
\[
 \boxed{m=\frac{qBB'\ell}{2E}}.
\]
同批同电荷态离子、相同的
\(E,B,B'\) 条件下质量与偏转
\(\ell\) 成正比：
\[
 \boxed{m_2=m_1\frac{\ell_2}{\ell_1}
 =16.00\,u\frac{32.86}{29.20}
 \approx18.01\,u},
\]
故第二种氧同位素为
\(\boxed{{}^{18}\mathrm O}\)。

\paragraph{习题 18.11}
载流电子数密度
\(n=8.4\times10^{22}\,\mathrm{cm^{-3}}
 =8.4\times10^{28}\,\mathrm{m^{-3}}\)。
霍尔平衡条件
\(|q|E_H=|q|v_dB\)，
电流密度
\(j=ne v_d\)。
片内电流通道横截面积
\(A=bd\)，
横向霍尔电压
\(U_H=E_Hb=v_dBb\)，
所以宽度 \(b\) 抵消：
\[
 \boxed{|U_{aa'}|=\frac{IB}{ned}
 =\frac{(200)(1.5)}
 {(8.4\times10^{28})(1.6022\times10^{-19})(.001)}
 \approx2.23\times10^{-5}\,\mathrm V}.
\]
即约 \(22.3\,\mu\mathrm V\)。
带符号电势 \(U_{aa'}\) 的正负
取决于图18.23中电流、电场箭头和
\(a,a'\) 位置；负载流子电荷
决定其与正载流子的霍尔符号相反。

\paragraph{习题 18.12}
图18.24中沿 \(x\) 的电流
\(I\)、沿 \(+z\) 的磁场，
霍尔场沿 \(y\)。
载流子数密度
\[
 n=\frac{IB}{|q|\,a\,|U_{AA'}|},
\]
因为电流横截面积为 \(ab\)，
而霍尔电压跨宽度 \(b\)，
故公式只剩厚度 \(a\)。
用
\(a=.001\,\mathrm m\)，
\(I=.001\,\mathrm A\)，
\(B=3000\,\mathrm G=.300\,\mathrm T\)，
\(U=6.55\,\mathrm{mV}\)，
\[
 \boxed{n\approx2.86\times10^{20}\,\mathrm{m^{-3}}}.
\]
\textbf{符号判定：}图示霍尔电势
\(U_{AA'}\) 必须连同图18.24
标示的A和A'两面实际位置一起审阅，
再以
\(\boldsymbol E_H=-\boldsymbol v_d\times\boldsymbol B\)
判定正负载流子。
首次候选在未严格分辨标注顺序前，
\emph{不将P或N型唯一化}；
数密度大小独立于此方向风险。

\paragraph{习题 18.13}
硅片长 \(\ell=.005\,\mathrm m\)，
宽 \(w=.002\,\mathrm m\)，
厚 \(d=.00005\,\mathrm m\)。
纵向通电电阻
\(R=\rho\ell/(wd)\)，
电流
\[
 I=\frac U R
 =\frac{Uwd}{\rho\ell}
 =\frac{(.002)(.00005)}{(.016)(.005)}
 \approx1.25\times10^{-3}\,\mathrm A.
\]
霍尔电压
\(U_H=IB/(ned)\)，
求磁场：
\[
 \boxed{B=\frac{U_H n e d}{I}
 =\frac{(.00105)(2\times10^{21})
 (1.6022\times10^{-19})(.00005)}{.00125}
 \approx1.35\times10^{-2}\,\mathrm T}.
\]
即约13.5mT；电阻率的
单位 \(\Omega\cdot\mathrm m\)
不可误看作 \(\Omega\)。

\paragraph{习题 18.14}
载流液体中体磁力密度
\(\boldsymbol f=\boldsymbol J\times\boldsymbol B\)。
当 \(\boldsymbol J\perp\boldsymbol B\)，
沿管道力密度大小 \(JB\)；
稳态流体沿管道有
\(dp/dx=JB\)，
在长度 \(l\) 上积分：
\[
 \boxed{\Delta p=JlB}.
\]
若 \(\Delta p=1.00\,\mathrm{atm}
 =101325\,\mathrm{Pa}\)，
\(B=1.50\,\mathrm T\)，
\(l=.0200\,\mathrm m\)，
\[
 \boxed{J=\frac{\Delta p}{lB}
 \approx3.38\times10^6\,\mathrm{A/m^2}}.
\]
这个理想模型忽略电极接触、电流不均匀、
壁摩擦和液体回流。

\paragraph{习题 18.15}
导电血液以速度 \(v\) 切割横向磁场时，
霍尔/动生电场 \(E=vB\)，
血管直径 \(d=.0020\,\mathrm m\)
两侧电压 \(U=E d\)，
\[
 \boxed{v=\frac U{Bd}
 =\frac{1.0\times10^{-4}}{(.080)(.0020)}
 \approx0.625\,\mathrm{m/s}}.
\]
实际交流信号的极性受探头方向影响，
此结果是速率大小。


\paragraph{习题 18.16}
图18.27的线圈一匝在磁场中的水平段
受竖直安培力大小 \(F_1=IlB\)，
\(n\) 匝为 \(nIlB\)。
电流反向时磁力反向，
天平两次状态的作用力\emph{差}为
\(2nIlB\)，因此所加砝码满足
\(mg=2nIlB\)。
\[
 \boxed{B=\frac{mg}{2nIl}
 =\frac{(.00878)(9.80)}
 {2(5)(.10)(.10)}
 \approx0.860\,\mathrm T}.
\]
不可遗漏电流反向所产生的因子2，
否则误报 \(1.72\,\mathrm T\)。

\paragraph{习题 18.17}
每匝面积
\(S=ab=(.020)(.010)=2.00\times10^{-4}\,\mathrm{m^2}\)，
匝数 \(N=1000\)，
\(I=.100\,\mathrm A\)，
\(B=1000\,\mathrm G=.100\,\mathrm T\)。
(a) 图18.28(a)磁场\emph{垂直}线圈平面：
对位于场内的每段直线
\[
 \boxed{\boldsymbol F_{\rm seg}
 =NI\boldsymbol\ell\times\boldsymbol B}.
\]
按图18.28(a)下水平边在场内，
其所受力大小
\[
 \boxed{F_{\rm lower}=NIaB
 =1000(.100)(.020)(.100)=0.200\,\mathrm N}.
\]
此边的力指向按原图电流箭头与入纸面磁场的
右手定则判断。
场中两条\emph{对称且等长}侧边所受横向力
大小相等、方向相反；
若上水平边完全在场外，则无该边磁力，
总平移力由下边决定。
由于线圈磁矩垂直线圈平面，
它与垂直外磁场平行或反平行，
\(\tau=0\)；
\emph{部分}线圈受场时还需以具体受力作用线
相对支点求转矩，不能无条件套整圈匀强场公式。

(b) 图18.28(b)磁场\emph{平行}线圈平面，
整个线圈均处于匀强场：
四边合力为零；
沿磁场方向的边有
\(d\boldsymbol\ell\parallel\boldsymbol B\)，
其力为零；
另两边分别受大小
\(F=NIaB=0.200\,\mathrm N\) 的等大反向力
（取图中与场垂直的边长为 \(a=.020\,\mathrm m\)），
形成力偶：
\[
 \boxed{\tau_{\max}=NISB
 =1000(.100)(2.00\times10^{-4})(.100)
 =2.00\times10^{-3}\,\mathrm{N\,m}}.
\]
这与磁矩
\(\boldsymbol\mu=NIS\,\hat{\boldsymbol n}\)
满足
\(\boldsymbol\tau=\boldsymbol\mu\times\boldsymbol B\)
吻合。
\textbf{审图风险：}上述边的“上/下/左右”及力箭头
以图18.28的原始电流箭头、支点为准，
后续TikZ重绘时必须逐边再次确认。

\paragraph{习题 18.18}
方形每边长 \(a=.150\,\mathrm m\)，
\(N=200\)，\(I=8.0\,\mathrm A\)，
面积 \(S=a^2=.0225\,\mathrm{m^2}\)。
线圈磁矩
\[
 \boxed{\mu=NIS=200(8.0)(.0225)
 =36.0\,\mathrm{A\,m^2}}.
\]
力矩 \(\tau=\mu B\sin\theta\)；
当磁矩方向与磁场垂直时最大，
\[
 \boxed{\tau_{\max}=\mu B=(36.0)(4.0)
 =144\,\mathrm{N\,m}}.
\]

\paragraph{习题 18.19}
圆盘面电荷密度 \(\sigma\)，
在半径 \(r\) 处宽 \(dr\) 的环上有
\(dq=2\pi r\sigma\,dr\)，
角速率 \(\omega\)，周期 \(2\pi/\omega\)，
等效环电流
\(dI=dq\,\omega/(2\pi)
 =\sigma\omega r\,dr\)。
该环磁矩 \(d\mu=dI\,\pi r^2\)，
所以
\[
 \boxed{\mu=\pi\sigma\omega\int_0^Rr^3dr
 =\frac14\pi\sigma\omega R^4}.
\]
总电量
\(q=\pi\sigma R^2\)，
盘的转动惯量
\(I_{\rm rot}=\frac12mR^2\)，
角动量
\(L=I_{\rm rot}\omega
 =\frac12mR^2\omega\)。
故
\[
 \boxed{\mu=\frac12q\omega R^2
 =\frac q{2m}L}.
\]
若 \(q<0\)，磁矩矢量与角动量反向，
大小式要用 \(|q|\)。

\paragraph{习题 18.20}
每个夸克沿半径 \(r\) 的圆轨道
速率大小 \(v\)，
相应角速率大小 \(\omega=v/r\)。
每个旋转点电荷的等效电流
\(I_i=q_i\omega_i/(2\pi)\)，
磁矩为
\(\mu_i=q_i\omega_i r^2/2\)（有符号）。
设两个下夸克沿正角向，
上夸克沿负角向，
\[
 \sum_iq_i\omega_i
 =2(-e/3)\omega+(+2e/3)(-\omega)
 =-\frac43e\omega.
\]
因此磁矩矢量沿负轴，大小
\[
 \boxed{|\mu|=\frac23e\omega r^2
 =\frac23evr}.
\]
按实验中子磁矩大小
\(\mu=9.66\times10^{-27}\,\mathrm{A\,m^2}\)
和原题半径
\(r=1.20\times10^{-15}\,\mathrm m\)，
\[
 \boxed{v=\frac{3\mu}{2er}
 \approx7.54\times10^7\,\mathrm{m/s}
 \approx0.251c}.
\]
该夸克经典绕行模型只用于估算，
真实中子磁矩涉及强相互作用与量子自旋，
不能据此认为夸克在精确圆周上运动。

\paragraph{习题 18.21}
若把电子自旋磁矩
\(\mu=0.928\times10^{-23}\,\mathrm{J/T}\)
误作单个电荷 \(e\) 以速率 \(c\)
绕半径 \(R\) 形成的经典环电流，
则
\[
 \mu=IA=(e c/(2\pi R))(\pi R^2)
 =\frac12ecR.
\]
于是
\[
 \boxed{R=\frac{2\mu}{ec}
 \approx3.86\times10^{-13}\,\mathrm m}.
\]
这个经典半径比实验给出电子结构尺度上界
\(10^{-18}\,\mathrm m\) 大约 \(4\times10^5\) 倍。
若强迫 \(R\le10^{-18}\,\mathrm m\)、
仍想保持相同磁矩，则所需表面运动速率
\(v=2\mu/(eR)\)
远大于 \(c\)，违反相对论。
结论：
\(\boxed{\text{电子自旋磁矩不是可旋转带电实心小球
的经典环流磁矩}}\)；
它是量子自旋本征性质。

\paragraph{习题 18.22}
图18.29中直导线电流 \(I\)，
线圈近侧电流 \(I_1\) 与它同向；
两者在同一平面。
距直导线 \(x\) 处
\(B(x)=\mu_0I/(2\pi x)\)，
线圈近侧长 \(l\)
受吸引力
\(F_{\rm near}=\mu_0II_1l/(2\pi a)\)。
远侧同长段电流反向，受排斥力
\(F_{\rm far}=\mu_0II_1l/[2\pi(a+b)]\)。
上下短边受大小相等、方向相反的磁力，
合力抵消，因此
\[
 \boxed{F_{\rm net}
 =\frac{\mu_0II_1l}{2\pi}
 \left(\frac1a-\frac1{a+b}\right)}
\]
方向\emph{朝向无限长直导线}。
线圈所在平面内的磁场始终垂直于线圈，
每个电流段磁力都位于该平面内；
刚性平面线圈在该对称分布下
总力矩为零：
\[
 \boxed{\tau_{\rm loop}=0}.
\]
若用户要求的“最大磁力矩”
指旋转线圈离开原图共面取向后的优化，
则需要额外定义旋转轴和近场几何，
不能沿用共面计算直接给非零数。

\paragraph{习题 18.23}
令无限薄圆筒半径为 \(R\)，
轴向面电流密度
\(\boldsymbol K=j\hat{\boldsymbol z}\)。
磁场在筒内为零，
筒外按安培定律
\[
 B_{\rm outside}(r)=\frac{\mu_0(2\pi Rj)}{2\pi r}
 =\frac{\mu_0jR}{r},
\]
在筒壁外侧 \(B_{\rm out}(R)=\mu_0j\)，
内侧 \(B_{\rm in}(R)=0\)。
自作用薄电流层受力必须取
\emph{两侧磁场平均}
\(\boldsymbol B_{\rm av}
 =\tfrac12(\boldsymbol B_{\rm out}+
 \boldsymbol B_{\rm in})\)，
不把自己电流元产生的奇异表面场重复算作外场。
单位面积磁力
\[
 \boxed{\boldsymbol f_A=
 \boldsymbol K\times\boldsymbol B_{\rm av}
 =-\frac{\mu_0j^2}{2}
 \hat{\boldsymbol r}}.
\]
大小为
\(\boxed{\mu_0j^2/2\,\mathrm{N/m^2}}\)，
方向\emph{指向圆筒轴线}（磁箍缩）。

\paragraph{习题 18.24}
把无限载流面看作
两侧磁场有跃变的理想电流薄片。
片面受磁力密度
\[
 \boxed{\boldsymbol f_A=
 \boldsymbol K\times
 \frac{\boldsymbol B_1+\boldsymbol B_2}{2}},
\]
这里两侧磁场应采用
\emph{矢量平均}，方向取原图18.30的
电流方向与磁场箭头确定。
对于两侧场皆切于片、垂直面电流、
并且平行/反平行于同一切向轴的情形，
边界条件
\(\mu_0\boldsymbol K
 =\hat{\boldsymbol n}\times
 (\boldsymbol B_2-\boldsymbol B_1)\)
给磁压的等价形式：
\[
 \boxed{|f_A|=\frac{|B_2^2-B_1^2|}{2\mu_0}}.
\]
方向由较高磁压侧指向较低磁压侧。
若图中的 \(B_1,B_2\)
是带方向的\emph{矢量}，不能把
\((B_1+B_2)/2\) 随意改成两正标量之和；
需要按右手定则和边界电流作矢量代数。

\paragraph{习题 18.25}
两条无限平行线距
\(d=.050\,\mathrm m\)，
各通过 \(I_1=I_2=30\,\mathrm A\)，
单位长度安培力
\[
 \boxed{\frac FL
 =\frac{\mu_0I_1I_2}{2\pi d}
 =\frac{(4\pi\times10^{-7})(30)^2}
 {2\pi(.050)}
 =3.60\times10^{-3}\,\mathrm{N/m}}.
\]
同向电流吸引，反向排斥。
当只有电子以缓慢平均漂移速率
\(10^{-3}\,\mathrm{m/s}\) 运输净电流时，
金属导线整体还含静止正离子晶格，
电荷总体近似中性，
\emph{不可}把全部移动电子之间的
库仑斥力当成两线总电力。
静电中性时跨导线的净库仑场强近似抵消，
而微小未抵消的速度相关场表现为磁相互作用。
完整相对论处理还与参考系及表面电荷分布有关。

\paragraph{习题 18.26}
半圆柱壳的总电流 \(I\)
沿轴向均匀分配在角度 \(\theta\in[0,\pi]\)，
微元电流
\(dI=(I/\pi)d\theta\)。
圆弧上各无限直电流丝
离轴心均为 \(R\)，
对轴心提供场微元大小
\(dB=\mu_0dI/(2\pi R)\)，
方向在截面内垂直半径。
沿半圆的垂直分量相互抵消，
沿对称轴的分量相加，
\[
 \boxed{B_{\rm axis}
 =\frac{\mu_0I}{2\pi^2R}
 \int_0^\pi\sin\theta\,d\theta
 =\frac{\mu_0I}{\pi^2R}}.
\]
与轴线电流 \(I\) 反向的直导线
所受单位长度磁力
\[
 \boxed{\frac FL
 =I B_{\rm axis}
 =\frac{\mu_0I^2}{\pi^2R}}.
\]
它受到半圆柱壳的\emph{排斥}，
方向沿壳面对称轴指向开口侧。
若用位于半圆壳对称轴上的
另一条同向于半圆壳的直导线
替代壳，轴上线与它反向电流的相互排斥力
大小为 \(\mu_0I^2/(2\pi d)\)。
由大小相等
\[
 \boxed{d=\frac{\pi R}{2}}.
\]
等效线放在\emph{原半圆柱面所在一侧}，
使它对中心轴线产生的磁力仍朝开口侧。
本题必须按图18.31区分开口侧与壳面侧，
不要把等效线放到相反一侧。

\paragraph{习题 18.27}
电磁轨道炮在题定平均匀加速模型中
\(v^2=2\bar aL\)。
\(v=4.2\times10^3\,\mathrm{m/s}\)，
\(L=5.0\,\mathrm m\)：
\[
 \boxed{\bar a=\frac{v^2}{2L}
 =1.764\times10^6\,\mathrm{m/s^2}
 =1.80\times10^5\,g}.
\]
按源书给定近似磁力公式
\[
 F=\frac{\mu_0I^2}{2\pi}
 \ln\frac{a+r}{r}=m\bar a.
\]
\(\,m=.317\,\mathrm{kg}\)，
\(a=.012\,\mathrm m,\ r=.067\,\mathrm m\)，
\(\ln[(a+r)/r]\approx.16476\)。
故
\[
 \boxed{I=\left[\frac{2\pi m\bar a}
 {\mu_0\ln((a+r)/r)}\right]^{1/2}
 \approx4.12\times10^6\,\mathrm A}.
\]
子弹动能
\(K=\frac12mv^2\approx2.796\times10^6\,\mathrm J\)；
匀加速历时
\(\Delta t=v/\bar a=2L/v\approx.002381\,\mathrm s\)。
若总效率40%且电源在发射期间近似
输出恒定平均功率，
\[
 \boxed{\bar P=\frac{K}{0.40\Delta t}
 \approx2.94\times10^9\,\mathrm W
 =2.94\times10^6\,\mathrm{kW}}.
\]
这是\emph{发射时段平均输入功率}，
不是电网日均耗电；
高脉冲功率通常需专用储能系统。
源数据 \(a<r\) 的几何解释仍应按原图18.32
复核，不擅改为不同尺寸。

\textbf{来源几何阻断项：}
现有候选题面称两圆柱轴距 \(a=1.2\,\mathrm{cm}\)，
半径 \(r=6.7\,\mathrm{cm}\)，
但分离圆柱的必要条件为 \(a\ge 2r\)，
实际 \(a/(2r)\simeq0.0896<1\)。
故上式 \(I\simeq4.12\,\mathrm{MA}\)、
\(\bar P\simeq2.94\,\mathrm{GW}\)
只能是代数公式的\emph{形式代入值}，
不能视为可实现导轨几何上的有效预测。
若原图中 \(a\) 另指\emph{净间隙}，或半径单位误录，
应先复核原书，再重新计算磁力积分；
不可自行猜改数值。


\paragraph{习题 18.28}
磁场垂直软导线所在平面，
长度 \(ds\) 的载流线元受安培力大小
\(dF=IB\,ds\) 沿圆弧法向。
曲线张力为 \(T\) 的细线
在半径 \(r\) 的小角元 \(d\theta\) 上，
两端张力的合向心力
\(2T\sin(d\theta/2)\simeq T\,d\theta\)。
由 \(ds=r\,d\theta\)：
\[
 T\,d\theta=IB\,r\,d\theta
 \Longrightarrow
 \boxed{r=\frac{T}{IB}}.
\]
另一方面带电量大小 \(q>0\)
的粒子作磁场圆运动，
\(qvB=mv^2/r\)，
故 \(r=p/(qB)\)。
令它与软导线几何圆半径相等，
\[
 \boxed{p=qBr=\frac{qT}{I}}.
\]
负电荷的绕行方向与正电荷相反，
但半径和动量\emph{大小}公式使用 \(|q|\)。

\paragraph{习题 18.29}
取相距矢量
\(\boldsymbol r=\boldsymbol r_2-\boldsymbol r_1
 =a\hat{\boldsymbol x}\)，
质子1速度沿连线远离质子2，
故 \(\boldsymbol v_1\parallel\boldsymbol r\)；
质子2速度
\(\boldsymbol v_2\perp\boldsymbol r\)，
为方便令 \(\boldsymbol v_2=v_2\hat y\)。
静电相互作用力大小
\(F_E=e^2/(4\pi\varepsilon_0a^2)\)，
两质子因同号而相互斥离。

低速带电点粒子产生的磁场近似
\[
 \boldsymbol B_1(\boldsymbol r_2)=
 \frac{\mu_0e}{4\pi a^2}
 \boldsymbol v_1\times\hat{\boldsymbol r}=\boldsymbol0
\]
（\(\boldsymbol v_1\parallel\hat{\boldsymbol r}\)），
所以质子2受质子1的磁力为零。
相反，质子2在质子1处的场
\(\boldsymbol B_2=
 \frac{\mu_0e}{4\pi a^2}
 \boldsymbol v_2\times(-\hat x)\)
一般不为零，
质子1受磁力
\[
 \boxed{\boldsymbol F_{1\leftarrow2}^{(B)}
 =\frac{\mu_0e^2}{4\pi a^2}
 \boldsymbol v_1\times[
 \boldsymbol v_2\times(-\hat x)]}
\]
方向按给定运动方向和右手法则求出。
在坐标约定
\(\boldsymbol v_1=-v_1\hat x\)
（质子1位于左、向左远离质子2）、
\(\boldsymbol v_2=+v_2\hat y\) 下，
磁力沿 \(+\hat y\)，其大小
\[
 \boxed{F_B=\frac{\mu_0e^2}{4\pi a^2}v_1v_2
 =F_E\frac{v_1v_2}{c^2}}.
\]
因此两质子\emph{机械}作用力
在此瞬间不严格互为相反矢量：
质子1有横向磁力而质子2没有。
这\emph{不}违反完整电磁系统动量守恒，
因为电磁场自身具有动量密度
\(\boldsymbol g=\varepsilon_0
 \boldsymbol E\times\boldsymbol B\)，
机械动量与场动量之和才守恒。
点电荷近场展开仅保留低速主项，
更完整的相对论场还包括迟滞项。

\paragraph{习题 18.30}
磁矩在外磁场中的势能
\(U=-\boldsymbol\mu\cdot\boldsymbol B
 =-\mu B\cos\theta\)。
\(\mu=2.39\times10^{-23}\,\mathrm{J/T}\)，
\(B=2.0\,\mathrm T\)，
\(\mu B=4.78\times10^{-23}\,\mathrm J\)。
将源书四种量子化方向代入，按角度顺序得到
\[
 \boxed{\begin{aligned}
 U(39.2^\circ)&\approx-3.70\times10^{-23}\,\mathrm J,\\
 U(75.0^\circ)&\approx-1.24\times10^{-23}\,\mathrm J,\\
 U(105.0^\circ)&\approx+1.24\times10^{-23}\,\mathrm J,\\
 U(140.8^\circ)&\approx+3.70\times10^{-23}\,\mathrm J.
 \end{aligned}}
\]
这对应原书磁矩相对于场方向的四种允许夹角；
以 \(U=0\) 选为垂直取向时的势能。
